\section{Hall correspondences on \texorpdfstring{$\mathfrak M^+_{\mathrm{eff}}(K3\times E)$}{the moduli of effective sheaves on K3 times E}}
\label{app:hall-correspondences}

The Hall side of the Dirac--Igusa construction rests on the following
imported theorems and on their \(K3\times E\) specialization.  The Hall
correspondences are the convolution correspondences through which the
Beilinson--Drinfeld factorization structure on the elliptic carrier
acts on the retained \(K3\times E\) Hall stack \cite{BD}.  Equivalently,
they comprise the chart-augmented \(A_b\) data for the open factorization dg-category
\((X, D, \tau)\) at a chosen boundary vacuum \(b\), with chiral shadow
on the elliptic carrier supplied separately by
\S~\ref{ch:zbps-realization}.

The cited theorems fix the exact hypotheses needed for the
\(K3\times E\) specialization: finite retained substacks, proper
extension correspondences, compact-support functoriality, vanishing
cycles, Thom--Sebastiani maps, and orientation coefficients.

The standing data are: \(S\) is a smooth complex projective \(K3\)
surface, \(E\) is a smooth elliptic curve, \(X = S\times E\) is the
trivial-canonical threefold, \(\sigma_S\in H^0(S, K_S)\) is the holomorphic
two-form trivializing \(K_S\), and \(\sigma = \mathrm{pr}_S^*\sigma_S
\in H^0(X, p_S^*K_S)\) is the inflated cosection of
\S\ref{ch:zbps-realization}.  The algebraic Mukai lattice used by the
Hall support is
\[
\Lambda_{S,\mathrm{alg}} = H^0(S,\Z)\oplus
\mathrm{NS}(S)\oplus H^4(S,\Z),
\]
the algebraic sublattice of the full Mukai lattice
\(\widetilde H(S,\Z)=H^0(S,\Z)\oplus H^2(S,\Z)\oplus H^4(S,\Z)\),
endowed with the restricted Mukai pairing of \cite{Mukai1984}.
\(\Gamma_{\mathrm{eff}}\) is the algebraic/effective Hall support of
Definition~\ref{def:algebraic-effective-hall-support}; and
\(\mathfrak M^+_{\mathrm{eff}}(X)\) is the moduli stack of perfect
complexes of effective dyonic class semistable for the supplied
stability datum on \(X\).

\subsection*{D.1\quad The Hall product on the moduli of stable sheaves}
\label{app:hall-product}

The first imported theorem is the existence and associativity of a
Hall product on the equivariant cohomology of the moduli stack of
semistable sheaves, with vanishing-cycle coefficients, after Joyce and
Song.

\begin{theorem}[Joyce--Song generalised \(\mathrm{DT}\) Hall product;
\(\gamma\), proved on \(\mathrm{CY}_3\) projective; conditional on
(O1), (O1)\(^+\) on \(K3\times E\)]
\label{thm:hall-product}
Let \(Y\) be a smooth complex projective Calabi--Yau threefold.  The
moduli stack \(\mathfrak M(Y)\) of compactly supported coherent
sheaves on \(Y\) carries a critical \emph{Hall product}
\[
m\;:\;
H^\bullet_c\bigl(\mathfrak M(Y),\Phi_W\otimes o\bigr)
\otimes
H^\bullet_c\bigl(\mathfrak M(Y),\Phi_W\otimes o\bigr)
\longrightarrow
H^\bullet_c\bigl(\mathfrak M(Y),\Phi_W\otimes o\bigr),
\]
with \(\Phi_W\) the perverse sheaf of vanishing cycles of the
holomorphic Chern--Simons potential and \(o\) the orientation line
supplied by Joyce--Upmeier orientation data.  The product is defined
by pull-push along the universal extension stack
\[
\mathfrak M(Y)\xleftarrow{\;p\;}
\mathfrak{Ext}^1(\mathfrak M(Y),\mathfrak M(Y))
\xrightarrow{\;q\;}
\mathfrak M(Y),
\]
where \(\mathfrak{Ext}^1\) is the moduli stack of short exact sequences
\(0\to A\to C\to B\to 0\), \(p(0\to A\to C\to B\to 0) = (A,B)\), and
\(q(0\to A\to C\to B\to 0) = C\).  The product
\(m = q_!\circ\mathrm{TS}\circ p^*\) is associative, where
\(\mathrm{TS}\) is the Massey--Saito--Joyce Thom--Sebastiani
isomorphism on the universal extension stack
\cite[Theorem~5.4]{JoyceSong}, \cite{MasseyST}, \cite{BBDJS2015}.  Its
unit is the class of the empty object.
\end{theorem}

\begin{proof}
The product is the pull-push form of three standard structures.  The
moduli stack \(\mathfrak M(Y)\) is locally the critical locus of the Joyce--Song
holomorphic Chern--Simons potential
\(W:\mathrm{Crit}(W)\hookrightarrow Y_W\) where \(Y_W\) is a Darboux
chart of the \((-1)\)-shifted symplectic structure on \(\mathfrak M(Y)\)
supplied by Pantev--To\"en--Vaqui\'e--Vezzosi
\cite[Theorem~0.3]{PTVV2013}; the perverse sheaf \(\Phi_W\) of
vanishing cycles is well defined on the oriented d-critical locus by
Brav--Bussi--Dupont--Joyce--Szendr\H{o}i
\cite[Theorem~6.9]{BBDJS2015}.  The universal extension stack
is locally a critical locus of \(W_A\boxplus W_B\) on
\(Y_A\times Y_B\), and the Massey--Saito Thom--Sebastiani isomorphism
\cite[Theorem~1.1]{MasseyST}, \cite[Theorem~5.10]{BBDJS2015} identifies
\(\Phi_{W_A\boxplus W_B}\) with \(\Phi_{W_A}\boxtimes\Phi_{W_B}\) on
\(\mathrm{Crit}(W_A)\times\mathrm{Crit}(W_B)\).  The
multiplicative orientation cocycle of Joyce--Upmeier
\cite[Theorem~4.1]{JoyceUpmeier},
\cite[Theorem~1.11]{JoyceTanakaUpmeier} trivializes the orientation
gerbe on extensions, supplying the line bundle factor.  When the
exceptional pushforward is admissible, pulling \(\Phi_W\otimes o\)
along \(p\), tensoring by Thom--Sebastiani, and pushing along \(q_!\)
defines \(m\).  Associativity reduces, via the
two-step flag stack
\(\mathfrak{Flag}(\mathfrak M(Y),\mathfrak M(Y),\mathfrak M(Y))\), to
the pentagon coherence of Thom--Sebastiani triple stratification
\cite[Theorem~5.10(iv)]{BBDJS2015} together with the Joyce--Upmeier
multiplicative cocycle on triple extensions
\cite[\S4]{JoyceUpmeier}.  See Joyce--Song
\cite[\S 5.2,~Theorem~5.5]{JoyceSong} for the original construction
and Davison \cite[Theorem~A]{Davison2017} for the cohomological
amplification.
\end{proof}

\noindent\emph{\(K3\times E\) specialization.} On \(Y = K3\times E\),
the cosection \(\sigma = \mathrm{pr}_S^*\sigma_S\) of the
holomorphic two-form on the K3 factor reduces the obstruction
theory through Kiem--Li (Appendix~\ref{app:cosection}); the supplied
moduli is \(\mathfrak M^{\mathrm{red}}_{\mathrm{eff}}(X)\subset
\mathfrak M(X)\), and the product is the reduced critical product
\[
m^{\mathrm{red}}\;=\;q_!^{\mathrm{red}}\circ\mathrm{TS}^{\mathrm{red}}
\circ p^*
\]
of Proposition~\ref{prop:finite-dcritical-hall-product}.  Vanishing of
the seven obstruction classes
\(\mathfrak o^{\mathrm{conf}}_R\),
\(\mathfrak o^{\mathrm{prop},\mathrm{LL}}_R\),
\(\mathfrak o^{\mathrm{prop},\mathrm{mix}}_R\),
\(\mathfrak o^{\mathrm{prop},\mathrm{WW}}_R\),
\(\mathfrak o^{\mathrm{conf\text{-}prop}}_R\),
\(\mathfrak o^{\mathrm{assoc}}_{w,R}\), and
\(\mathfrak o^{I,\mathrm{prod}}_{w,R}\), together with finite residual
inertia, retained compactified flag-Quot correspondences or an
explicit compact-support six-functor model, base-change,
projection-formula, and quotient-after-correspondence coherences, is
the data \((D0)\)+\((O1)\), absorbed into \(\mathrm{Bnd}\),
\(\mathrm{dCrit}\), \(\mathrm{RedOr}\), and \(\mathrm{Six}\) of
Definition~\ref{def:compact-hall-product}.  \emph{The
existence of a chain-level Hall product on \(K3\times E\) without
those four data is open.} The resulting invariant
\(Y^+(X) := H^\bullet_c(\mathfrak M^+_{\mathrm{eff}}(X), \Phi_W\otimes o)\)
is the positive-half cohomological Hall algebra (CoHA) of
Kontsevich--Soibelman \cite{KSCoHA} restricted to effective dyonic
class; its Drinfeld double, regulated tensor product completion, and
integral form constitute the level-3 object \(G(X)\), conditional on
the Hall pairing and stable-envelope transport.  Thus \(G(X)\) arises
from \(Y^+(X)\) by Hall pairing, integral-form completion, and
stable-envelope transport, as in the Vol~III \(\Delta_5\) row
\cite{CYQGVolIII}.

\subsection*{D.2\quad The Hall coproduct via collisions}
\label{app:hall-coproduct}

The Hall product does not by itself supply a coproduct.  A collision
coproduct is additional six-functor data on the reversed extension
correspondence: properness or compact-support replacement for the
left leg, inverse Thom--Sebastiani transport, quotient descent, vacuum
unit/counit data, and the four-step flag-of-flags compatibility that
compares \(\Delta m\) with
\((m\otimes m)(1\otimes\tau\otimes1)(\Delta\otimes\Delta)\).  The
Hopf pairing and its radical quotient are separate data.  When the
coproduct and bialgebra compatibility are supplied, \(Y^+(X)\) becomes
a graded bialgebra object; after the Hopf pairing, Drinfeld doubling,
completion, integral form, and stable-envelope transport, it is a
candidate for the level-3 quantum group \(G(X)=D_\hbar(Y^+(X))\) of
Schiffmann--Vasserot type \cite{SchiffmannVasserot2013}.

\begin{theorem}[Collision coproduct datum, formal conditional]
\label{thm:hall-coproduct}
Let \(Y\) be a smooth projective Calabi--Yau threefold and suppose
the Joyce--Song Hall product of Theorem~\ref{thm:hall-product} is
given.  Assume, in addition, that the reversed extension
correspondence carries admissible compact-support six-functor data,
inverse Thom--Sebastiani transport, vacuum unit/counit data, and
three-step and four-step flag coherences for coassociativity and
bialgebra compatibility.  Then the critical Hall positive half
\(Y^+(Y) = H^\bullet_c(\mathfrak M^+_{\mathrm{eff}}(Y),\Phi_W\otimes o)\)
carries a collision coproduct
\[
\Delta\;:\;Y^+(Y)\longrightarrow Y^+(Y)\otimes Y^+(Y),
\qquad
\Delta\;=\;p_!^{\mathrm{coll}}\circ
\mathrm{TS}^{-1}\circ q^{\mathrm{coll},*},
\]
defined by the collision/extension correspondence
\[
\mathfrak M^+_{\mathrm{eff}}(Y)
\xleftarrow{\;q^{\mathrm{coll}}\;}
\mathfrak{Ext}^{1,\mathrm{coll}}(\mathfrak M^+_{\mathrm{eff}}(Y))
\xrightarrow{\;p^{\mathrm{coll}}\;}
\mathfrak M^+_{\mathrm{eff}}(Y)\times\mathfrak M^+_{\mathrm{eff}}(Y),
\]
with \(\mathfrak{Ext}^{1,\mathrm{coll}}\) the moduli of short exact
sequences \(0\to A\to C\to B\to 0\) read with \(C\) on the left and
\((A,B)\) on the right, \(q^{\mathrm{coll}}(0\to A\to C\to B\to 0) = C\),
\(p^{\mathrm{coll}}(0\to A\to C\to B\to 0) = (A,B)\).  The
coproduct \(\Delta\) is graded coassociative.  Together with the Hall
product \(m\) of Theorem~\ref{thm:hall-product},
\((Y^+(Y), m, \Delta)\) is a graded bialgebra.
\end{theorem}

\begin{proof}
The formula is the formal reverse of the Hall product.  Its existence
requires the extra six-functor hypotheses because the reversed leg is
not automatically proper in the compact threefold problem.
Given those hypotheses, coassociativity is the inverse
Thom--Sebastiani pentagon on the three-step flag stack.  The
compatibility of \(m\) and \(\Delta\) is not a consequence of the
product theorem alone; it is the equality of the two pull-push functors
on the four-step flag-of-flags common refinement.  On quivers with
potential this is the Davison--Meinhardt bialgebra theorem
\cite[Theorem~A]{DavisonMeinhardtCoHA}; on the compact
\(K3\times E\) threefold it is the finite correspondence datum of
Definition~\ref{def:compact-geometric-hall-bialgebra-datum} and
Proposition~\ref{prop:geometric-hall-data-give-bialgebra-defects}.
\end{proof}

\noindent\emph{Drinfeld double.} The Hall pairing
\(\langle\,\cdot\,,\,\cdot\,\rangle_{\mathrm{Hall}}\), when supplied
and non-degenerate after the completed radical quotient, pairs
\(Y^+\) with itself; the Drinfeld double \(D_\hbar(Y^+(X))\) is the
quantum double in the sense of Drinfeld \cite[\S 13]{KAC:1}, supplying
the negative half \(Y^-(X)\) and the Cartan generators \(K_\gamma\),
\(\gamma\in\Gamma_{\mathrm{eff}}\), with relations encoding the Hall
pairing.  The \emph{level-3 \(K3\times E\) quantum group} is then
\[
G(X)\;:=\;D_\hbar(Y^+(X))_{\mathrm{compl}}^{\mathrm{int}}
\]
where \((-)_{\mathrm{compl}}^{\mathrm{int}}\) is the integral-form
completion under the Hall pairing, conditional on the
stable-envelope transport of \S~D.6.  In Schiffmann--Vasserot
\cite{SchiffmannVasserot2013} the analogous identification on
toric \(\mathrm{CY}_3\) ambient is established;
\(K3\times E\) is not equipped here with a global toric or
quiver-variety NCCR model of that kind: there is no finite
fixed-point McKay presentation, no compatible product-polarized HPD
self-dual model, and no Schiffmann--Vasserot stable-envelope atlas on
the compact non-toric ambient.  The cosection-twisted reduced
obstruction theory of \S~D.4 is the replacement; the global
\(G(X)\) on the compact non-toric \(K3\times E\) is not constructed
from these data without the Hall pairing, integral form, completion,
and stable-envelope transport.

\subsection*{D.3\quad The Hall primitives and the formal recognition criterion}
\label{app:hall-primitives}

For a connected graded bialgebra \(B = \bigoplus_{c}B_c\), the
primitive elements at charge \(c\) are
\[
\mathrm{Prim}_c(B)\;=\;
\{x\in B_c\mid \Delta(x) = x\otimes 1 + 1\otimes x\},
\]
equivalently the kernel of the reduced coproduct \(\bar\Delta\).  On a
graded connected cocommutative Hopf algebra in characteristic zero,
the Milnor--Moore theorem \cite[Theorem~5.18]{LurieHA} identifies
\(B\simeq U(\mathrm{Prim}_\bullet(B))\) as Hopf algebras, where
\(\mathrm{Prim}_\bullet(B)\) is the Lie algebra of primitives under
the bracket \([x,y] = xy - (-1)^{|x|\,|y|}yx\).

\begin{theorem}[Hall primitives, formal level; \(\beta\), formal]
\label{thm:hall-primitives-formal}
Let \(\mathfrak M\) carry a critical Hall positive half
\(Y^+(Y) = (B, m, \Delta)\) of
Theorem~\ref{thm:hall-coproduct}.  The graded subspace
\[
\mathrm{Prim}_\bullet(Y^+(Y))\;=\;
\bigoplus_{c\in\Gamma_{\mathrm{eff}}}
\mathrm{Prim}_c(Y^+(Y))
\]
is a graded Lie subalgebra of \(Y^+(Y)\) under the antisymmetrised
product \([\,\cdot\,,\,\cdot\,]\).  At a charge \(c\) such that the
moduli \(\mathfrak M^+_c(Y)\) is connected and the \(C\)-component of
\(\Delta\) on \(\mathfrak M^+_c(Y)\) factors through extensions whose
graded pieces lie in strictly lower charges, the inclusion
\(\mathrm{Prim}_c\hookrightarrow Y^+_c\) is the kernel of
\(\bar\Delta_c\).  When \(Y^+(Y)\) is graded connected and
cocommutative, Milnor--Moore identifies
\(Y^+(Y)\simeq U(\mathrm{Prim}_\bullet(Y^+(Y)))\) as graded Hopf
algebras.
\end{theorem}

\begin{proof}
Primitivity of \(\mathrm{Prim}_\bullet\) under the antisymmetrised
product is verified by direct computation of
\(\Delta([x,y])\) using the cocommutative-up-to-twist axiom and the
multiplicativity of \(\Delta\); the resulting expression collapses to
\([x,y]\otimes 1 + 1\otimes [x,y]\).  On a connected graded
cocommutative Hopf algebra in characteristic zero, the Milnor--Moore
identification is \cite[\S 21.1]{LurieHA}; on graded connected
non-cocommutative Hopf algebras, the Cartier--Quillen--Milnor--Moore
theorem \cite[Theorem~21.2.1]{LurieHA} requires either
cocommutativity or a chain-level lift through bar / cobar, which is
exactly the chiral Koszul separation theorem of Vol I
\cite{FrancisGaitsgoryCKD}.
\end{proof}

\noindent\emph{Recognition.}
The recognition theorem
of Theorem~\ref{thm:primitive-recognition} identifies the completed
formal primitive Lie superalgebra with the imported
Borcherds--Kac--Moody Lie superalgebra
\(\mathfrak g_{\Delta_5}\) of
Theorem~\ref{thm:bkm-denominator-algebra-identity}, conditional on the
finite compact-source datum \((S1)\)--\((S10)\): the completed PBW
filtration, the Hopf pairing radical equality, and the parity
dimension recovery.  In
particular, on \(K3\times E\) the formal primitive subspace
\(\mathrm{Prim}_\bullet(Y^+(X))\) is the \emph{candidate} primitive
Lie superalgebra; its identification with the imported BKM
denominator algebra \(\mathfrak g_{\Delta_5}\) is the primitive
recognition part of View \(\textcircled{3}\).  The formal Hopf algebra
structure alone gives no recognition theorem; recognition acts on it.

\subsection*{D.4\quad The cosection-twist on \texorpdfstring{$K3\times E$}{K3 x E}: Kiem--Li localisation}
\label{app:cosection}

On \(K3\times E\), the holomorphic two-form on the K3 factor supplies
the semiregularity cosection on the retained nonzero K3-class
branches of the obstruction theory of \(\mathfrak M(X)\).  The
Kiem--Li \cite{KiemLi2013} localisation theorem then localises the
virtual fundamental class to the zero-locus of the cosection when the
surjectivity hypothesis holds.  The localisation is the source of the
\emph{reduced} virtual class on which the Hall product of
\S~D.1 carries chain-level coherence.

\begin{definition}[Cosection]
\label{def:cosection-app-D}
Let \(\mathfrak M\) be a Deligne--Mumford stack equipped with a
perfect obstruction theory \(\phi:\mathbb E\to\mathbb L_{\mathfrak M}\)
of amplitude \([-1,0]\), and let \(\mathrm{Ob}(\mathfrak M) =
h^1(\mathbb E^\vee)\) be the coherent obstruction sheaf.  A
\emph{cosection} of the obstruction theory is an
\(\mathcal O_{\mathfrak M}\)-linear morphism
\[
\sigma\;:\;\mathrm{Ob}(\mathfrak M)\longrightarrow
\mathcal O_{\mathfrak M};
\]
its \emph{vanishing locus} is the closed substack
\(\mathfrak M^{\mathrm{red}} = (\sigma = 0)\) on which the obstruction
sheaf restricts to the kernel \(\ker\sigma\subset\mathrm{Ob}(\mathfrak M)\).
\end{definition}

\begin{theorem}[Kiem--Li cosection localisation;
\cite{KiemLi2013}~Theorem~5.0.1; \(\gamma\), proved]
\label{thm:cosection-localisation}
Let \(\mathfrak M\) be a Deligne--Mumford stack with perfect
obstruction theory \(\phi\) and surjective cosection
\(\sigma:\mathrm{Ob}(\mathfrak M)\twoheadrightarrow\mathcal O_{\mathfrak M}\)
on \(\mathfrak M\setminus\mathfrak M^{\mathrm{red}}\), where
\(\mathfrak M^{\mathrm{red}} = (\sigma = 0)\).  The virtual
fundamental class \([\mathfrak M]^{\mathrm{vir}}\in
A_{\mathrm{vd}}(\mathfrak M)\) lifts to a cosection-localised class
\[
[\mathfrak M]^{\mathrm{vir}}_{\mathrm{loc}}
\;\in\;A_{\mathrm{vd}}(\mathfrak M^{\mathrm{red}})
\]
satisfying \(\iota_*[\mathfrak M]^{\mathrm{vir}}_{\mathrm{loc}} =
[\mathfrak M]^{\mathrm{vir}}\), where
\(\iota:\mathfrak M^{\mathrm{red}}\hookrightarrow\mathfrak M\) is the
inclusion.  If the cosection is used separately to quotient a trivial
obstruction direction and form a reduced obstruction theory, that
reduced theory has virtual dimension
\(\mathrm{vd}^{\mathrm{red}}=\mathrm{vd}+1\); this dimension shift is
not the Kiem--Li localized cycle itself.
\end{theorem}

\begin{proof}
The cosection-kernel cone \(C(\sigma) = \ker(\sigma:\mathrm{Ob}\to
\mathcal O)\) is a perfect obstruction cone of amplitude \([-1,0]\)
on \(\mathfrak M^{\mathrm{red}}\) with
\(\mathrm{rk}\,C(\sigma) = \mathrm{rk}\,\mathrm{Ob} - 1\).  The
intrinsic normal cone \(\mathfrak C\) of \(\mathfrak M\) sits in the
total obstruction bundle on \(\mathfrak M\setminus\mathfrak M^{\mathrm{red}}\),
where \(\sigma\) is surjective and the cone has codimension one in
\(\mathrm{Ob}\); on \(\mathfrak M^{\mathrm{red}}\), the cone embeds
into \(C(\sigma)\) by \cite[Proposition~4.3]{KiemLi2013}.
Intersection with the zero section of \(C(\sigma)\) on
\(\mathfrak M^{\mathrm{red}}\), pushed forward to \(\mathfrak M\)
through \(\iota\), yields \([\mathfrak M]^{\mathrm{vir}}\) by
\cite[Theorem~5.0.1]{KiemLi2013}.  The localized Gysin map preserves
the original virtual dimension.  The shift
\(\mathrm{vd}^{\mathrm{red}} = \mathrm{vd} + 1\) enters only after a
separate reduced obstruction theory is formed by removing the trivial
obstruction summand.
\end{proof}

\noindent\emph{\(K3\times E\) cosection.} On \(X = K3\times E\), the
cosection of the obstruction theory of \(\mathfrak M(X)\) is induced
by the pullback \(\sigma=\mathrm{pr}_S^*\sigma_S\) of the K3
holomorphic two-form.
Concretely, for a perfect complex \(F\) on \(X\), the
chain-level cosection acts on \(\mathrm{Ext}^2(F,F)\) by
the semiregularity map of Buchweitz--Flenner /
Bloch \cite{MaulikTodaGV},
\[
\mathrm{cs}_F\;:\;\mathrm{Ext}^2(F,F)\longrightarrow
\C,
\qquad
\mathrm{cs}_F(\xi)\;=\;\int_X
\sigma\wedge\mathrm{Tr}\bigl((\xi\otimes 1_{\Omega_X^1})\circ
\mathrm{At}(F)\bigr),
\]
where \(\mathrm{At}(F)\in\mathrm{Ext}^1(F, F\otimes\Omega^1_X)\) is
the Atiyah class.  The cosection vanishes precisely on the
\emph{reduced} moduli
\(\mathfrak M^{\mathrm{red}}(X)\subset\mathfrak M(X)\) consisting of
sheaves whose Atiyah class is annihilated by \(\sigma\); equivalently,
sheaves whose deformation space along the K3 direction is
constrained.  See Maulik--Toda
\cite[Lemma~2.6, Theorem~2.7]{MaulikTodaGV} and Oberdieck
\cite[Theorem~1, Theorem~3]{OberdieckReducedPT} for the explicit
identification on the DT/PT side.  Surjectivity is an additional
row-level condition: the retained stratum must be identified with the
nonzero K3-class branch and must record a cokernel-rank-zero
semiregularity row.  On the \(s=0\) Hall cusp the retained charge can
lose the K3 component that makes the semiregularity map surjective; the
holomorphic two-form \(\sigma_S\) does not vanish.  The localisation then passes to the
\(D0\)-degeneration limit named at obstruction \((D0)\) of
\S~\ref{ch:zbps-realization}.

\begin{theorem}[Reduced virtual class on \(K3\times E\); primitive
stable-pair branch]
\label{thm:reduced-virtual-class-K3xE}
On the stable-pair or ideal-sheaf Deligne--Mumford branch carrying a
perfect obstruction theory and a surjective K3 semiregularity
cosection, Kiem--Li localization gives
\[
[\mathfrak M_c(X)]^{\mathrm{vir}}_{\mathrm{loc}}\in
A_{\mathrm{vd}}(\mathfrak M_c(X)^{\mathrm{red}}).
\]
The reduced obstruction theory obtained by quotienting the trivial
obstruction direction gives a distinct reduced class
\[
[\mathfrak M_c(X)]^{\mathrm{red}}\in
A_{\mathrm{vd}+1}(\mathfrak M_c(X)^{\mathrm{red}}).
\]
The reduced DT/PT theory of the primitive nonzero K3-class branch uses
this reduced class and gives the Oberdieck--Pixton scalar
\(Z^X_{\mathrm{OP}}=-4096\,\Delta_5^{-2}\)
\cite[Theorem~3]{OberdieckReducedPT}.  The semistable Hall stack
requires the retained compact-support six-functor and d-critical Hall
datum of \S\ref{subsec:hall-correspondences}.
On the \(s = 0\) cusp the cosection degenerates; \emph{the
cosection-reduced theory there is the \((D0)\) degeneration limit
named at \S~\ref{ch:zbps-realization} and controlled by the retained
D0-HN criterion of
Theorem~\ref{thm:G-D0}}.\quad[\(\delta\), \((D0)\)]
\end{theorem}

\subsection*{D.5\quad BBDJS reduced d-critical orientation}
\label{app:bbdjs}

The Hall product of \S~D.1 requires not just a virtual fundamental
class but a perverse-sheaf-of-vanishing-cycles refinement and an
orientation gerbe; the source is the Brav--Bussi--Dupont--Joyce--Szendr\H{o}i
\cite{BBDJS2015} extension of the Joyce
\cite{JoyceSong} d-critical-locus framework.

\begin{theorem}[BBDJS perverse-sheaf-of-vanishing-cycles; oriented
\(\mathrm{CY}_3\) version of \cite{BBDJS2015}~Theorem~6.9; \(\gamma\), proved]
\label{thm:bbdjs}
Let \(\mathfrak M\) be a derived Artin stack with a \((-1)\)-shifted
symplectic structure (PTVV \cite{PTVV2013}) whose classical truncation
is a d-critical locus in the sense of Joyce
\cite[Definition~2.5]{JoyceSong}.  Suppose \(\mathfrak M\) carries an
\emph{orientation}, i.e.~a square root \(o\) of the Joyce determinant
line \(K_{\mathfrak M}^{\mathrm{vir}} = (\det\mathbb L_{\mathfrak M})^{1/2}\),
in the sense of \cite[Theorem~6.9]{BBDJS2015}.  There is a perverse
sheaf
\[
\Phi_{\mathfrak M}\;\in\;\mathrm{Perv}(\mathfrak M),
\]
the \emph{perverse sheaf of vanishing cycles} of the d-critical
locus; locally on a Joyce--Song Darboux chart \(\mathrm{Crit}(W)\)
inside a smooth ambient \(Y_W\), \(\Phi_{\mathfrak M}|_{\mathrm{Crit}(W)}
\simeq\Phi_W[d_{Y_W}]\) is the standard perverse sheaf of vanishing
cycles of the holomorphic potential \(W:Y_W\to\C\), shifted by the
ambient dimension.  The chart-comparison cocycle is canonical up to
the orientation choice; the orientation determines a multiplicative
cocycle of Joyce--Upmeier \cite[Theorem~4.1]{JoyceUpmeier} on the
moduli of objects.
\end{theorem}

\begin{proof}
Three layers stack.  PTVV \cite[Theorem~0.3]{PTVV2013} produces the
\((-1)\)-shifted symplectic form on the derived moduli of
\(\mathrm{CY}_3\) coherent sheaves; Bussi--Brav--Joyce
\cite[\S 5]{BBDJS2015} produce the d-critical-locus chart structure
on the classical truncation; Bussi--Brav--Dupont--Joyce--Szendr\H{o}i
\cite[Theorem~6.9]{BBDJS2015} produce the perverse sheaf of
vanishing cycles with its chart cocycle.  The orientation is the datum needed
to unify chart-local \(\Phi_W\) into a global perverse sheaf
\(\Phi_{\mathfrak M}\); without it, the chart-cocycle has a sign
ambiguity that obstructs the descent to a perverse sheaf
\cite[\S 6.4]{BBDJS2015}.  Joyce--Upmeier
\cite[Theorem~1.11]{JoyceUpmeier} construct the multiplicative orientation
cocycle on the universal extension stack; this is the
\(\mathrm{CY}_3\) projective theorem for the retained Hall orientation.
\end{proof}

\noindent\emph{Reduced refinement on \(K3\times E\).}
On \(X = K3\times E\), the \((-1)\)-shifted symplectic structure
\(\omega_{\mathfrak M}\) on \(\mathfrak M(X)\) does not restrict to a
\((-1)\)-shifted symplectic structure on
\(\mathfrak M^{\mathrm{red}}(X)\); the cosection-twist of \S~D.4
contracts \(\omega_{\mathfrak M}\) by \(\sigma\), and the
\emph{reduced} \((-1)\)-shifted form
\(\omega_{\mathfrak M}^{\mathrm{red}}\) is non-degenerate only after
restriction to the cosection vanishing locus \cite[\S 2.6]{MaulikTodaGV}.
The corresponding reduced d-critical-locus structure is supplied by
the Maulik--Toda \cite[Theorem~2.7]{MaulikTodaGV} extension; the
reduced perverse sheaf of vanishing cycles
\(\Phi_{\mathfrak M}^{\mathrm{red}}\) and the reduced orientation
line \(o^{\mathrm{red}}_{\mathcal C}\) are exactly the data
\((O1)\) of \S~\ref{ch:zbps-realization}.  In
the notation of (P4), the reduced
orientation is twisted by a parity-flip line \(\Pi\) generated by the
cosection cokernel,
\[
o^{\mathrm{red}}_{\mathcal C}\;=\;
o^{\mathrm{std}}_{\mathcal C}\otimes\Pi,
\]
absorbing the dimension shift \(\mathrm{vd}^{\mathrm{red}} =
\mathrm{vd}^{\mathrm{std}} + 1\) into a \(\Z/2\)-grading shift on the
oriented BBDJS line.

\begin{proposition}[BBDJS reduced d-critical orientation on
\(K3\times E\); \(\delta\), \((O1)+(O1)^+\)]
\label{prop:bbdjs-reduced-K3xE}
Let \(\mathfrak M^{\mathrm{red},+}_c(X)\) be the cosection-reduced
moduli of semistable sheaves of effective dyonic class \(c\) on
\(X = K3\times E\).  The hypotheses \emph{strong reduced orientation
on the \(K3\times E\)-quotient self-Ext frame bundle (O1) and
Weyl-equivariant transport along \(W^{(2)}(\Lambda_{II}^{2,1})\)
reflections (O1)\(^+\)} of
\S~\ref{ch:zbps-realization} are exactly the
following.  There exist a multiplicative orientation cocycle of
Joyce--Upmeier type extending equivariantly across the reduced
\(E\)-quotient, the \(\sigma_S\)-cosection-induced parity flip line
\(\Pi\), and a connected \(BE\)-equivariant trivialisation of the
small-orbit gerbe class \(\overline\beta^E_{\mathcal C}\) on the
finite-stabilizer strata, such that
\[
\Phi^{\mathrm{red}}_{\mathfrak M^{\mathrm{red},+}_c(X)}
\otimes
o^{\mathrm{red}}_c
\]
descends through the free \(E\)-translation quotient and trivializes
the Pin\(^c\)-lift gerbe class of the reduced derived self-Ext frame
bundle.\quad The strong reduced orientation \((O1)\) is established
from the retained quotient-orientation datum.  Joyce--Upmeier
\cite[Theorem~1.11]{JoyceUpmeier} supplies the projective
Calabi--Yau-threefold orientation technology used here; the
quasi-projective and quotient descent belongs to the retained
\((O1_{\mathrm{quot}})\) datum.
The Cao--Kool--Monavari CY\(_4\) theorem
\cite[Theorem~1.6]{CaoKoolMonavari} is a fourfold analogue, not an
ingredient of the \(K3\times E\) CY\(_3\) reduction.  The free
\(E\)-quotient, finite-stabilizer, and Weyl-transport theorem on
\(K3\times E\) follows only relative to the retained
quotient-orientation datum and the chosen Weyl lifts of
Theorems~\ref{thm:G-O1} and~\ref{thm:G-O1-plus}.
\end{proposition}

\begin{proof}
The construction of the reduced perverse sheaf
\(\Phi^{\mathrm{red}}_{\mathfrak M^{\mathrm{red},+}_c(X)}\) on the
cosection-reduced moduli is by Maulik--Toda
\cite[Theorem~2.7]{MaulikTodaGV}; the orientation gerbe lifts to the
reduced Joyce d-critical locus by the BBDJS
\cite[\S 6.4]{BBDJS2015} chart-cocycle calculation.  The
multiplicative \emph{equivariant} orientation cocycle across the
\(E\)-quotient and across the finite-stabilizer locus
\(BE[N]\)-strata is the joint vanishing of \((O1)\) and \((O1)^+\):
on the connected \(BE\)-equivariant trivialisation of
\(\alpha^E_{\mathcal C}\in H^2_E(\mathfrak M^{\mathrm{red}}_{\mathcal C};\mathbb F_2)\)
of Definition~\ref{def:O1-plus-weyl-transport} it suffices to trivialize the
small-orbit gerbe class \(\overline\beta^E_{\mathcal C}\) on
\(BE[N]\)-strata, governed by the Klein-four
detection criterion of Lemma~\ref{lem:G-Klein-four}.  The Weyl-equivariant transport across
\(W^{(2)}(\Lambda_{II}^{2,1})\) reflections is the orientation
character calculation of
Theorem~\ref{thm:pfaffian-dirac-finite-stage}
on the three type-II walls.  Theorems~\ref{thm:G-O1}
and~\ref{thm:G-O1-plus} supply these transports on \(K3\times E\);
the Pin\(^c\)-lift theory of
\cite{JoyceUpmeier,JoyceTanakaUpmeier} supplies the analytic
technology in the Calabi--Yau-threefold setting.
\end{proof}

\subsection*{D.6\quad Stable-envelope transport: from Hall to Yangian}
\label{app:stable-envelope}

The level-3 quantum-group structure on \(G(X) = D_\hbar(Y^+(X))\) of
\S~D.2 is, on toric ambient or quiver varieties, identified by
Maulik--Okounkov \cite{MaulikOkounkov} with the geometric \(R\)-matrix
of the \emph{stable envelope}, providing the link from the Hall side
to the Yangian / quantum-group side.

\begin{definition}[Stable envelope]
\label{def:stable-envelope}
Let \(\mathfrak M\) be a smooth symplectic variety with a torus
\(T\)-action and a finite \(T\)-fixed locus \(\mathfrak M^T\); let
\(\mathcal C\subset\mathrm{Lie}(T)\) be a chamber of the
\(T\)-equivariant cohomology of \(\mathfrak M\).  The
\emph{stable envelope} \cite[\S 3]{MaulikOkounkov} is the
\(T\)-equivariant cohomological correspondence
\[
\mathrm{Stab}_{\mathcal C}\;:\;
H^\bullet_T(\mathfrak M^T)\longrightarrow
H^\bullet_T(\mathfrak M),
\]
characterised by support, normalization, and degree axioms
\cite[Theorem~3.3.4]{MaulikOkounkov}.  The geometric \(R\)-matrix is
the comparison
\[
R^{\mathrm{MO}}_{\mathcal C, \mathcal C'}\;=\;
\mathrm{Stab}_{\mathcal C'}^{-1}\circ\mathrm{Stab}_{\mathcal C}
\;\in\;\mathrm{End}(H^\bullet_T(\mathfrak M^T))
\]
across two chambers \(\mathcal C, \mathcal C'\); it satisfies the
Yang--Baxter equation by \cite[Theorem~4.6.1]{MaulikOkounkov}.
\end{definition}

\begin{theorem}[Maulik--Okounkov; \cite{MaulikOkounkov}~Theorem~10.2.1;
\(\beta\), proved on quiver varieties; conditional on \(K3\times E\)
compact non-toric]
\label{thm:maulik-okounkov}
Let \(\mathfrak M = \mathfrak M_{Q,\mathbf v,\mathbf w}\) be a Nakajima
quiver variety \cite{Nakajima1994,Nakajima1998} attached to a finite
quiver \(Q\) with dimension vectors \(\mathbf v, \mathbf w\).  The
stable envelope \(\mathrm{Stab}_{\mathcal C}\) makes
\(H^\bullet_T(\mathfrak M_{Q,\mathbf v,\mathbf w})\) a module over a
Yangian \(Y_\hbar(\mathfrak g_Q)\) attached to \(Q\), with
\(\mathfrak g_Q\) the Kac--Moody Lie algebra of the quiver, and the
geometric \(R\)-matrix \(R^{\mathrm{MO}}\) is the rational
\(R\)-matrix of \(Y_\hbar(\mathfrak g_Q)\).
\end{theorem}

\begin{proof}
The Maulik--Okounkov construction proceeds through:
(i) the existence and uniqueness of the stable envelope (axioms of
support, polarisation, normalization), \cite[\S 3.3]{MaulikOkounkov};
(ii) the Yang--Baxter equation as a chamber-wall cocycle,
\cite[\S 4]{MaulikOkounkov}; (iii) the identification of the Bethe
algebra generated by the matrix elements of \(\mathrm{Stab}\) with a
Yangian, \cite[\S 5--10]{MaulikOkounkov}.  The Schiffmann--Vasserot
identification \(\mathrm{CoHA}(\C^3) = Y^+(\widehat{\mathfrak{gl}}_1)\)
of \cite[Theorem~A]{SchiffmannVasserot2013} is the affine \(A_0\)
specialization; \(\mathcal W_{1+\infty}\) arises \emph{only after}
Drinfeld doubling and Fock-evaluation.  The three-arrow chain
\(\mathrm{CoHA} = Y^+\hookrightarrow Y\xrightarrow{\mathrm{ev}_\lambda}
\mathrm{End}(\mathcal W_{1+\infty}[\lambda]\text{-vac})\) of
\cite{KSCoHA}, \cite{SchiffmannVasserot2013} traverses the level-2
\(\to\) level-3 \(\to\) level-3-evaluated passage; the named identity
on each arrow measures the structural distance from
\(\mathrm{CoHA}(\C^3)\) to \(\mathcal W_{1+\infty}\).
\end{proof}

\noindent\emph{\(K3\times E\) specialization: open.}
On \(X = K3\times E\), the moduli \(\mathfrak M^+_{\mathrm{eff}}(X)\)
is not a Nakajima quiver variety, the ambient is compact non-toric,
and the global stable envelope is unconstructed.  The five
obstructions to a global toric/quiver-variety NCCR model of the
Schiffmann--Vasserot type obstruct the literal
\cite{MaulikOkounkov} construction.  The cosection-twist of \S~D.4 supplies the
\emph{Serre-equivariant reduced substitute}: the
cosection-reduced obstruction theory of
Theorem~\ref{thm:reduced-virtual-class-K3xE} restores a
\(\mathrm{CY}_3\)-symmetric reduced obstruction theory on the
\(s\ne 0\) sector, and the Hall-to-Yangian comparison is then
expected at the level of the reduced theory and conditional on the
stable-envelope transport, after the retained D0-HN datum, retained
quotient-orientation datum, chosen Weyl lifts, and retained
\((O2_{\mathrm{atlas}})\) wall datum of
Appendix~\ref{app:obstruction-discharges}.

The Hall--Yangian comparison on \(K3\times E\) is therefore
conjectural:
\[
G(X)\;\overset{?}{=}\;
D_\hbar(\mathrm{CoHA}^{\mathrm{red},+}_{K3\times E})_{\mathrm{compl}}^{\mathrm{int}}
\;\overset{?}{=}\;
Y_\hbar(\widehat{\mathfrak g}_{\Delta_5})
\]
where \(\widehat{\mathfrak g}_{\Delta_5}\) is the central extension of
the imported BKM Lie superalgebra of
Theorem~\ref{thm:bkm-denominator-algebra-identity}, and the equality is
conditional on a \(K3\times E\) stable-envelope transport.  The
Hall--Drinfeld double \(\mathcal D_\hbar(\mathrm{CoHA}_{K3\times E})\)
of Kontsevich--Soibelman \cite{KSCoHA} places this identification at
the same level-3 quantum group seen from the geometric side; it is
the open compact non-toric instance of the chain-fusion conjecture.

\subsection*{D.7\quad Level separations}
\label{app:hall-level-separations}

The Beilinson cut separates the following objects.

\begin{enumerate}
\item \(Y^+(X)\) is the cohomological-Hall positive half
\(H^\bullet_c(\mathfrak M^+_{\mathrm{eff}}(X), \Phi_W\otimes o)\) with
the Joyce--Song product/coproduct of
Theorems~\ref{thm:hall-product},~\ref{thm:hall-coproduct}; \(G(X) =
D_\hbar(Y^+(X))_{\mathrm{compl}}^{\mathrm{int}}\) is the Drinfeld
double after Hall pairing, completion, integral form, stable-envelope
transport, and descent.

\item \(\mathrm{CoHA}(\C^3)=Y^+(\widehat{\mathfrak{gl}}_1)\)
\cite[Theorem~A]{SchiffmannVasserot2013}; \(\mathcal W_{1+\infty}\)
appears only after Drinfeld doubling and Fock evaluation.

\item No global Schiffmann--Vasserot/Nakajima atlas is known for compact non-toric
\(K3\times E\).  The cosection-twist of \S~D.4 supplies the
Serre-equivariant reduced substitute used here.

\item \(\mathrm{Prim}_\bullet(Y^+(X))\) is the formal primitive Lie
superalgebra of \(Y^+(X)\); its identification with
\(\mathfrak g_{\Delta_5}\) is the primitive-recognition theorem of
Theorem~\ref{thm:primitive-recognition}, conditional on the finite
compact-source datum \((S1)\)--\((S10)\).

\item BBDJS \cite[Theorem~6.9]{BBDJS2015} requires an
\emph{oriented} d-critical locus; the orientation is the square root
of the Joyce determinant line, and the multiplicative cocycle is
Joyce--Upmeier \cite[Theorem~4.1]{JoyceUpmeier}.

\item Kiem--Li \cite[Theorem~5.0.1]{KiemLi2013} requires the cosection to be
surjective off the localised locus; on \(K3\times E\), the cosection
degenerates on the \(s = 0\) cusp, which is the \((D0)\)
degeneration limit controlled by the retained D0-HN criterion of
Theorem~\ref{thm:G-D0}.

\item The cosection contracts the \((-1)\)-shifted symplectic form;
\(\omega^{\mathrm{red}}\) is non-degenerate only after
restriction to the cosection vanishing locus
\cite[\S 2.6]{MaulikTodaGV}, and the orientation line carries the
\(\Z/2\)-parity flip \(\Pi\) absorbing the dimension shift
\(\mathrm{vd}^{\mathrm{red}} = \mathrm{vd}^{\mathrm{std}} + 1\).

\item \(K3\times E\) is a trivial-canonical threefold.  The scalar-trace
index is \(n=\chi(\mathcal O_Y)\) on the
closed target, whereas
\(\chi(\mathcal O_{K3\times E})=\chi(\mathcal O_{K3})\chi(\mathcal O_E)
=2\cdot0=0\), and the BBDJS / Joyce--Upmeier
machinery applies in the \(\mathrm{CY}_3\) form
\cite[\S 6]{BBDJS2015}, not in the \(\mathrm{CY}_4\) Borisov--Joyce
\cite{BorisovJoyce} or Oh--Thomas \cite{OhThomasI,OhThomasII} form.
\end{enumerate}

\subsection*{D.8\quad Four local obstruction calculations}
\label{app:hall-local-obstruction-calculations}

Four local calculations enter the Hall correspondence.

\textbf{(F1)\quad Hall associativity.} The Joyce--Song associativity
of Theorem~\ref{thm:hall-product} requires the pentagon coherence of
Thom--Sebastiani triple stratification
\cite[Theorem~5.10(iv)]{BBDJS2015}.  The
chart-cocycle of the d-critical locus carries the orientation sign
whose pentagon coherence is measured by the BBDJS orientation gerbe.
That gerbe \cite[\S 6.4]{BBDJS2015} trivializes
exactly this ambiguity; the multiplicative cocycle of Joyce--Upmeier
\cite[\S 4]{JoyceUpmeier} extends the trivialisation to the universal
extension stack.  On \(K3\times E\), the gerbe descends through the
free \(E\)-quotient and trivializes the small-orbit class by
Theorems~\ref{thm:G-O1} and~\ref{thm:G-O1-plus}.

\textbf{(F2)\quad Cosection surjectivity.} Kiem--Li
\cite[Theorem~5.0.1]{KiemLi2013} requires
\(\sigma:\mathrm{Ob}\twoheadrightarrow\mathcal O\) on
\(\mathfrak M\setminus\mathfrak M^{\mathrm{red}}\).  On the \(s = 0\)
cusp of the dyonic charge lattice of \(K3\times E\),
the K3-class component needed for semiregularity can disappear and
surjectivity fails, while the K3 holomorphic two-form itself remains
nonzero.  The smooth primitive nonzero K3-class branch satisfies
surjectivity by Maulik--Toda
\cite[\S 2]{MaulikTodaGV}; the cusp branch is the \((D0)\)
degeneration limit named at obstruction \((D0)\) of
\S~\ref{ch:zbps-realization}, controlled by the retained D0-HN
criterion of Theorem~\ref{thm:G-D0}.
The Oberdieck--Pixton scalar
\(Z^X_{\mathrm{OP}} = -4096\,\Delta_5^{-2}\) is proved on the
\(s\ne 0\) branch \cite[Theorem~3]{OberdieckReducedPT}; on the cusp
the chain-level theory degenerates to the Borcherds
imaginary-simple-root contribution.

\textbf{(F3)\quad BBDJS orientation existence on compact non-toric.}
The Joyce--Upmeier multiplicative cocycle \cite[Theorem~4.1]{JoyceUpmeier}
is established on the \(\mathrm{CY}_3\) projective ambient; its
extension across the cosection-reduced quotient ambient is part of the
retained \((O1_{\mathrm{quot}})\) datum.  On the \(K3\times E\)
quotient by free \(E\)-translation, descent of the orientation line
requires the small-orbit gerbe class on \(BE[N]\)-strata to
trivialize.  The joint vanishing of the connected
\(BE\)-equivariant degree-two obstruction
\(\alpha^E_{\mathcal C}\in H^2_E(\mathfrak M^{\mathrm{red}}_{\mathcal C};\mathbb F_2)\)
the finite-stabilizer Cech--Borel gerbe classes, and the residual
degree-one stabilizer characters is exactly \((O1)\);
Theorem~\ref{thm:G-O1} of Appendix~\ref{app:obstruction-discharges}
is the quotient-orientation criterion relative to this retained datum.
The remaining open data here are the retained quotient orientation and
the compact-source recognition rows.

\textbf{(F4)\quad \((-1)\)-shifted vs \((-2)\)-shifted symplectic.}
\(K3\times E\) is a trivial-canonical threefold; the moduli
\(\mathfrak M(K3\times E)\) carries a \((-1)\)-shifted symplectic
structure by PTVV \cite[Theorem~0.3]{PTVV2013}.  The Borisov--Joyce
\cite{BorisovJoyce} and Oh--Thomas \cite{OhThomasI,OhThomasII}
theories are \(\mathrm{CY}_4\) theories.  The Calabi--Yau dimension
of \(X = K3\times E\) is \(\dim_{\C}(K3) + \dim_{\C}(E) = 2+1 = 3\),
the holomorphic three-form is \(\sigma_S\wedge dz_E\), and the
shifted symplectic is \([-1]\)-shifted by \cite[\S 2.1]{PTVV2013}.
The cosection of \S~D.4 removes one obstruction direction and
increases the reduced virtual dimension by one, but does not reduce
the shifted-symplectic degree; \(\omega^{\mathrm{red}}\) is still
\((-1)\)-shifted, on the cosection vanishing locus, by the
Maulik--Toda \cite[Theorem~2.7]{MaulikTodaGV} extension.  The
cosection-twist of \S~D.4 is the Serre-equivariant reduced substitute
for the missing global toric/quiver-variety atlas.

The four calculations above are the local Hall forms of
\((D0)\), \((O1)\), \((O1)^+\), and \((O2)\).  Their vanishing is the
hypothesis of the conditional cross-level identity
\(\textcircled{2}\)\(\leftrightarrow\)\(\textcircled{3}\).  The Hall
correspondences above supply the compact-support functorial maps; the
recognition
\(\mathrm{Prim}_\bullet(Y^+(X))\simeq\mathfrak g_{\Delta_5}\) of
Theorem~\ref{thm:primitive-recognition} remains conditional on
\((S1)\)--\((S10)\).  The protected Pfaffian theorem compares the
View~\textcircled{3} Pfaffian section with \(\mathcal D_X=\Delta_5\); the
level-\(\mathsf Z\) scalar trace is a later trace operation.
